% Бүлэг 1

\chapter{Судалгааны хэсэг} % Бүлгийн нэр
\label{Chapter1} % Энэ бүлэг рүү ишлэл хийх бол \ref{Chapter1} командыг ашигла 

%-------------------------------------------------------------------------------


%-------------------------------------------------------------------------------
\section*{Системийн зорилго}

Пиццаны захиалгыг хүлээн авахаас эхлээд хүргэлт хийх, төлбөр хүлээн авах, захиалгыг хаах хүртэлх үйл ажиллагааг нэгдсэн системээр бүртгэн, хянаж, хурдан шуурхай, үнэн зөв зохион байгуулахад оршино.


\section{Системийн зорилт}
\begin{enumerate}
    \item Утсаар ирсэн захиалгыг компьютерт хурдан, зөв бүртгэх.
    \item Захиалагчийн мэдээлэл болон идэвхтэй захиалга байгаа эсэхийг шалгах.
    \item Идэвхтэй захиалгатай хэрэглэгчээс давхар захиалга авахгүй байх.
    \item Захиалгын нийт үнийг автоматаар тооцоолох.
    \item Нийт үнэ 40,000 төгрөгөөс дээш бол 20\%-ийн хөнгөлөлтийг автоматаар олгох.
    \item Захиалгыг хүргэгчид хуваарилж, хүргэлтийн баримт хэвлэх.
    \item Захиалгын бүртгэл, бэлтгэл болон хүргэлтийн төлөвийг хянах.
    \item Хүргэгчийн хүлээн авсан төлбөр, гарын үсэгтэй баримтын буцаалтыг бүртгэх.
    \item Мөнгө болон баримтыг шалгасны дараа захиалгыг хаах.
    \item Борлуулалт, хөнгөлөлт, төлбөр болон хүргэлтийн тайлан гаргах.
    \item Алдаа, мэдээллийн зөрүү болон төлбөрийн эрсдэлийг бууруулах.
    \item Захиалгын үйлчилгээний хурд, хяналт, найдвартай байдлыг сайжруулах.
\end{enumerate}


%-------------------------------------------------------------------------------
\section{Программын судалгаа}
\subsection{Дотоод системтэй харьцуулсан судалгаа}

\subsubsection{Pizza Hut Mongolia систем}

Манай тодорхойлсон систем нь утсаар захиалга хүлээн авч, хүргэлт болон бэлэн мөнгөний тооцоог хянах дотоод систем юм. Харин Pizza Hut Mongolia нь вэбсайт, гар утасны апп, хэрэглэгчийн бүртгэл бүхий илүү өргөн хүрээний онлайн захиалгын систем ашигладаг. Албан ёсны сайтаар хүргэлт эсвэл салбараас авах захиалга хийх, бүтээгдэхүүн болон багц сонгох боломжтой.

\textbf{www.pizzahut.mn-ийн гол онцлогууд:}
\begin{table}
    \centering
    \begin{tabular}{p{4cm}p{4.5cm}p{5.5cm}}
        \toprule
\textbf{Үзүүлэлт} &
\textbf{Таны систем} &
\textbf{Pizza Hut Mongolia} \\
\midrule

Захиалгын суваг &
Зөвхөн утсаар &
Вэбсайт, апп, утас \\

Захиалга бүртгэх &
Ажилтан компьютерт гараар бүртгэнэ &
Хэрэглэгч өөрөө онлайн захиалах боломжтой \\

Хэрэглэгчийн бүртгэл &
Утасны дугаараар шалгана &
Нэр, утас, имэйл, нууц үгээр бүртгэл үүсгэнэ \\

Бүтээгдэхүүн сонгох &
Захиалга бүртгэгч сонгоно &
Хэрэглэгч цэс, багц, пицца, нэмэлт, ундаанаас сонгоно \\

Бүтээгдэхүүн тохируулах &
Тодорхойлоогүй &
Пиццаны хэмжээ, гурилын төрлийг тохируулж болно \\
Хөнгөлөлт &
40,000-өөс дээш бол тогтмол 20\% &
Багц, урамшуулал болон купон код ашиглана \\

Захиалгын давхардал &
Идэвхтэй захиалгатай хэрэглэгчээс дахин авахгүй &
Нийтийн сайтаас ийм хязгаарлалт илрээгүй \\

Хүргэлтийн төрөл &
Зөвхөн хүргэлт &
Хүргэлт эсвэл салбараас авах \\

Хүргэгчид шилжүүлэх &
Хэвлэсэн баримтаар &
Дотоод удирдлагын системтэй байх шинжтэй боловч нарийн ажиллагаа
нийтэд нээлттэй биш \\
Төлбөр &
Хүргэгч бэлэн мөнгө хураана &
Төлбөрийн сонголт захиалгын хэсэгт хийгдэнэ; нийтэд харагдах
мэдээллээр бэлэн мөнгөний сонголт байна \\

Баталгаажуулалт &
Баримтад гарын үсэг зуруулна &
Цахим захиалгын мэдээлэл хадгалагдана \\

Захиалга хаах &
Мөнгө, баримт тушаасны дараа ажилтан хаана &
Дотоод процесс нь нийтэд дэлгэрэнгүй харагдахгүй \\

Тайлан, удирдлага &
Тайлангийн модуль шаардлагатай &
Тусдаа менежер болон ERP нэвтрэх системүүдтэй \\

\bottomrule

    \end{tabular}
    \caption{Пицца захиалгын системүүдийн харьцуулалт}
    \label{tab:placeholder}
\end{table}

Pizza Hut Mongolia-ийн апп нь цэс харах, пиццаны сонголтыг тохируулах, захиалга өгөх боломжтой. Мөн сайт нь хэрэглэгчийн нэвтрэх, бүртгүүлэх болон купон код ашиглах хэсэгтэй.
Уг хийж буй системийг Pizza Hut-ийн жишигт ойртуулахын тулд дараах боломжийг нэмэх нь тохиромжтой:

\begin{itemize}
    \item Вэб болон гар утсаар захиалга өгөх.
    \item Хүргэлт эсвэл салбараас авах сонголт.
    \item Цэс, үнэ, бүтээгдэхүүний үлдэгдлийн удирдлага.
    \item Пиццаны хэмжээ, орц, гурилыг тохируулах.
    \item Купон болон олон төрлийн урамшууллын дүрэм.
    \item QPay, банкны карт зэрэг цахим төлбөр.
    \item Захиалгын төлөвийг хэрэглэгчид харуулах.
    \item Хүргэгчийн байршил, хүрэх хугацааг хянах.
    \item SMS эсвэл аппын мэдэгдэл илгээх.
    \item Давтан захиалга болон дуртай бүтээгдэхүүн хадгалах.
\end{itemize}
Ингэснээр анхны систем зөвхөн дотоод бүртгэлийн хэрэгсэл биш, хэрэглэгч–ажилтан–гал тогоо–хүргэгч–төлбөрийг холбосон нэгдсэн захиалга, хүргэлтийн систем болно.

\section{Хууль, дүрэм журам}
Монгол Улсад онлайн пицца захиалгыг дангаар зохицуулсан тусгай хууль одоогоор байхгүй. Харин цахим худалдаа, хэрэглэгчийн эрх, хувийн мэдээлэл, татвар болон хүнсний аюулгүй байдлын хэд хэдэн хууль хамт үйлчилнэ.


\subsection{Хамаарах хууль, заалт}
\begin{table}
    \centering
    \begin{tabular}{p{4cm}p{4.5cm}p{5.5cm}}
\toprule
\textbf{Хууль, эрх зүйн акт} &
\textbf{Гол заалт} &
\textbf{Системд тавигдах шаардлага} \\
\midrule

\textbf{Хэрэглэгчийн эрхийг хамгаалах тухай хууль} &
4.1.2, 5.1, 7.1--7.2 &
Бүтээгдэхүүний нэр, орц, найрлага, хэмжээ, үнэ, хадгалах болон
хэрэглэх мэдээллийг үнэн зөв харуулна. \\

\textbf{Иргэний хууль} &
39--43, 196 дугаар зүйл &
Захиалга баталгаажсан мөч, худалдагч ба худалдан авагчийн үүрэг,
цуцлалтын нөхцөлийг тодорхой болгоно. \\

\textbf{Хүний хувийн мэдээлэл хамгаалах тухай хууль} &
8, 18, 20, 22 дугаар зүйл &
Нэр, утас, хаяг, байршил зэрэг мэдээллийг зөвшөөрөлтэй цуглуулж,
хамгаална. \\

\textbf{Хүнсний бүтээгдэхүүний аюулгүй байдлыг хангах тухай хууль} &
3.1, 9, 10 дугаар зүйл &
Пиццаны орц, түүхий эдийн гарал, хадгалалт, үйлдвэрлэл, хүргэлтийн
бүртгэл хөтөлнө. \\

\textbf{Хүнсний тухай хууль} &
10 дугаар зүйл болон холбогдох заалт &
Хүнсний чиглэлийн үйл ажиллагаа эрхлэгч аюулгүй, шаардлага хангасан
хүнс нийлүүлнэ. \\

\textbf{Нэмэгдсэн өртгийн албан татварын тухай хууль} &
17.3--17.5 &
Борлуулалт бүрийг бүртгэж, төлбөрийн баримт буюу eBarimt олгоно. \\

\textbf{Цахим гарын үсгийн тухай хууль} &
3 дугаар зүйл &
Албан ёсны цахим баримт бичиг, гэрээнд шаардлагатай бол цахим эсвэл
тоон гарын үсэг хэрэглэнэ. \\

\textbf{Хоол үйлдвэрлэл, үйлчилгээний газрын эрүүл ахуй, халдвар
хамгааллын үлгэрчилсэн дүрэм} &
1.3--1.7 болон холбогдох бүлгүүд &
Гал тогоо, тоног төхөөрөмж, ажилтны ариун цэвэр, хадгалалт,
цэвэрлэгээ, дотоод хяналтыг зохицуулна. \\

\bottomrule
        
    \end{tabular}
    \caption{Хуулийн судалгаа}
    \label{tab:placeholder}
\end{table}

Хэрэглэгч барааны талаар бодит мэдээлэл авах эрхтэй бөгөөд үйлдвэрлэгч нэр, хаяг, бүтээгдэхүүний орц, найрлага, нэр төрөл, үнэ, хэмжээ зэрэг мэдээллийг өгөх үүрэгтэй.

\subsection{Хувийн мэдээллийн шаардлага}
Онлайн систем дараах мэдээллийг цуглуулж болно:
\begin{itemize}
    \item захиалагчийн нэр;
    \item утасны дугаар;
    \item хүргэлтийн хаяг;
    \item байршлын мэдээлэл;
    \item захиалгын түүх;
    \item төлбөрийн төлөв.
\end{itemize}

Мэдээлэл цуглуулахаас өмнө хэрэглэгчид зорилго, цуглуулах мэдээллийн жагсаалт, хадгалах хугацаа, бусдад дамжуулах эсэх болон зөвшөөрлөө цуцлах аргыг мэдээлнэ. Зөвшөөрлийг цаасан эсвэл цахим хэлбэрээр авч болно.

Системд дараах хамгаалалт шаардлагатай:
\begin{itemize}
    \item хэрэглэгчийн зөвшөөрлийг бүртгэх;
    \item нууцлалын бодлого харуулах;
    \item ажилтны эрхийг албан тушаалаар хязгаарлах;
    \item хүргэгчид зөвхөн шаардлагатай нэр, утас, хаягийг харуулах;
    \item нууц үгийг хамгаалалттай хадгалах;
    \item мэдээллийн өөрчлөлтөд аудитын лог үүсгэх;
    \item мэдээлэл алдагдсан үед авах арга хэмжээний төлөвлөгөөтэй байх;
    \item хэрэглэгч мэдээллээ засах, устгуулах хүсэлт гаргах боломжтой байх.
\end{itemize}

\subsection{Хүнсний аюулгүй байдал, хүргэлт}

Пицца бэлтгэхээс хүргэх хүртэл хүнсний аюулгүй байдлыг хянах хэрэгтэй. Хууль нь хадгалах, тээвэрлэх, худалдах, үйлчлэх бүх үе шатанд үйлчилнэ.

Системд дараах бүртгэл байвал зохино:
\begin{itemize}
    \item орц, түүхий эдийн нийлүүлэгч;
    \item хүлээн авсан огноо болон цуврал;
    \item хадгалалтын хугацаа, температур;
    \item пицца бэлтгэсэн цаг;
    \item захиалга хүргэлтэд гарсан цаг;
    \item хүргэгчийн мэдээлэл;
    \item хүргэж дууссан цаг;
    \item гомдол, бүтээгдэхүүн буцаалт;
    \item аюултай бүтээгдэхүүнийг татан авсан бүртгэл.
\end{itemize}

Хүнсний үйл ажиллагаа эрхлэгч нь түүхий эдийн ул мөрийг мөрдөн тогтоох, дотоод хяналтын баримт хадгалах, ажилтнаа эрүүл ахуйн сургалт болон эрүүл мэндийн үзлэгт хамруулах, бүтээгдэхүүнийг эрүүл ахуйн шаардлагын дагуу хадгалж тээвэрлэх үүрэгтэй.

\subsection{Төлбөр ба eBarimt}

НӨАТ суутган төлөгч бол систем:
\begin{itemize}
    \item Төлбөр хийсэн тухай бүр борлуулалтыг бүртгэнэ.
    \item Бэлэн мөнгө, карт, QPay зэрэг төлбөрийн төрлийг хадгална.
    \item Хөнгөлөлтийн өмнөх болон дараахах үнийг ялгаж харуулна.
    \item Худалдаа бүрд eBarimt үүсгэнэ.
    \item Баримтын мэдээллийг борлуулалтын нэгдсэн системд хуульд заасан хугацаанд илгээнэ.
    \item Цуцлагдсан болон буцаагдсан төлбөрийг тусад нь бүртгэнэ.
\end{itemize}

НӨАТ-ын тухай хуулийн 17.3-т борлуулалт бүрд төлбөрийн баримт олгох, мэдээллийг гурван хоногийн дотор нэгдсэн системд илгээхээр заасан


\subsection{Хамаарах стандартууд}

\begin{table}
    \centering
    \begin{tabular}{p{5cm}p{10cm}}
        \toprule
\textbf{Стандарт} &
\textbf{Хэрэглэх хүрээ} \\
\midrule

\textbf{MNS 4946:2019} --- Хоол үйлдвэрлэл, үйлчилгээний газарт
тавих шаардлага &
Пиццаны газар, гал тогоо, тоног төхөөрөмж, үйлчилгээ, ажилтны нөхцөл \\

\textbf{MNS ISO 22000:2019} --- Хүнсний аюулгүй байдлын менежментийн
тогтолцоо &
Хүнсний эрсдэлийн удирдлага, HACCP, ул мөр, хяналтын тогтолцоо \\

\textbf{MNS CAC RCP 1} --- Хүнсний эрүүл ахуйн ерөнхий зарчим &
Түүхий эдээс хэрэглэгч хүртэлх эрүүл ахуйн зохистой дадал \\

\textbf{MNS 6648:2016} --- Хүнсний сав, баглаа боодлын шошгололтод
тавих шаардлага &
Савласан ундаа, соус болон бусад савласан бүтээгдэхүүний шошго \\

\textbf{MNS 7032:2024} --- Түргэн хөргөх, хөлдөөсөн хоол, бэлэн
бүтээгдэхүүний шаардлага &
Хөлдөөсөн пицца эсвэл урьдчилан бэлтгэсэн бүтээгдэхүүн борлуулах үед \\

\bottomrule
    \end{tabular}
    \caption{Caption}
    \label{tab:placeholder}
\end{table}

MNS 4946:2019 нь хоол үйлдвэрлэл, үйлчилгээний газрын үндсэн салбарын стандарт юм. MNS ISO 22000:2019 нь хүнсний аюулгүй байдлын удирдлага болон HACCP тогтолцоонд тавих шаардлагыг тодорхойлдог

\subsection{Системд заавал тусгах мэдээлэл}

Захиалга баталгаажуулахын өмнө хэрэглэгчид дараахыг харуулна:

\begin{enumerate}
    \item байгууллагын нэр, хаяг, холбоо барих мэдээлэл;
    \item пиццаны нэр, хэмжээ, орц, харшил үүсгэгч;
    \item үндсэн үнэ;
    \item хөнгөлөлтийн хувь ба нөхцөл;
    \item хүргэлтийн төлбөр;
    \item төлөх нийт дүн;
    \item төлбөрийн хэлбэр;
    \item хүргэлтийн хаяг ба тооцоолсон хугацаа;
    \item захиалга цуцлах, буцаалт хийх нөхцөл;
    \item гомдол гаргах суваг;
    \item нууцлалын бодлого.
\end{enumerate}

Анхаарах нь: Үндэсний стандарт бүр автоматаар заавал мөрдөх шинжтэй биш. Харин хууль, техникийн зохицуулалт, тусгай дүрэм, гэрээ эсвэл зөвшөөрлийн нөхцөлд эш татсан бол заавал мөрдөх шаардлага үүсэж болно. Иймээс үйл ажиллагаа явуулах газрын дүүрэг, хүнсний үйлдвэрлэлийн төрөл болон борлуулах бүтээгдэхүүнээр нь нэмэлт шаардлагыг дахин шалгах хэрэгтэй. Энэ жагсаалтыг 2026 оны 9 дүгээр сарын 22-ны байдлаар албан эх сурвалжид тулгуурлан нэгтгэв.
\section{Нэр томъёоны тайлбар}
	\begin{longtable}{|p{0.5cm}|p{4cm}|p{8.5cm}|}
            \hline
                № & Тодорхойлолт & Тайлбар \\ 
            \hline
                1 & Task Management & Төсөл хуваарилах системийн нэр  \\ 
            \hline
               
        \caption{Нэр томъёоны тайлбар}
        \label{table:1}
    \end{longtable}